\section{Consequences for a learner that can revise its representation}
\label{sec:yba-discussion}

The register at the beginning could move, arbitrate or split. We can now say what distinguishes those operations. Register motion reduces a task score at a fixed authority. Arbitration redistributes whose error can move the same substrate. Splitting changes which demands must share a state. The common ledger supplies an exact boundary at which the last operation pays. The boundary would be invisible to a model that allowed only weight changes.

An endogenous evaluator introduces a further state. Our explicit reuse-mode controller shows that its authority coordinate can converge together with the model under stated conditions. It does not implement or establish convergence of a capacity-purchase mechanism. The purchase boundary remains a static comparison, and the tracking policies use their own fixed scoring rule. The damping, coupling, parent drive and update order carry the reuse-mode result. A switched architecture with a learning parent must include those additional states in its own analysis.

The tracking evidence then separates coherent accounting from a successful evaluator. Actual-step pricing assigns the executed move its actual observed-window gain and propagates its state correctly. Yet it predicts worse than the noisy-step comparator at every matched price in the retained experiment. The productive conclusion is not that cost accounting should be abandoned. It is that the evaluator's horizon and statistical assumptions remain design questions. A policy can reliably buy what it scored while scoring the wrong thing for tomorrow's prediction.

The representation calculation suggests how to make those questions precise. OR aggregation preserves the event exactly if the receiver retains cumulative hazard. Summed margins with one common intercept lose a particular distinction between blocks. Their residual TV measures disagreement about the event. Repricing asks a richer question and needs more than one hazard value. Adaptive evolution asks a richer question again and requires a state or history kernel. Each enlargement is motivated by an operation the smaller representation cannot support.

This gives a practical interpretation of the proposed three laws. Endogenous evaluation asks where the score comes from and how it changes. Embodied price asks which transition is purchased, on which ledger and horizon. Representation covariance asks which map, observable or transition is to survive, at what tolerance. A complete answer need not be exact at every scale. It must make the approximation and the receiving question explicit.

Several comparisons now become possible without invoking a universal theory of intelligence. Hold task losses fixed and compare a weight-only controller with one that can purchase an additional register at a declared price. Hold executable actions fixed and compare observed-window gain with a forecast-based or uncertainty-sensitive evaluator. Hold fine dynamics fixed and compare receiving descriptions with and without retained history. In each case the extra operation has an explicit cost and an explicit consequence to measure.

The nested construction also gives a disciplined use for attention-like arbitration and finite self-modelling. Authority weights can route error among competing demands without creating capacity. A represented self-state can help a controller predict its own subsequent operation without requiring unlimited recursion. Whether either mechanism pays in a particular organism or engineered system is a question about its actual signals, resources and task. The finite framework supplies candidate measurements, not a biological guarantee or a clinical diagnosis.

The most useful next step is therefore not to announce that all learning minimises one hidden scalar. It is to build models in which the alternatives are genuine: a move that is executed as priced, an evaluator that can be perturbed, a partition whose purchase changes the feasible states, and a receiving representation whose predictions can be compared. Self-consistent loss geometry names the coupled object that remains when these choices are made explicit. Its scientific value lies in the calculations and failures those choices let us expose.
